\section{Survey-Derived Edge-Contract Reference Framework}
\label{sec:synthesis}

The five gaps in \cref{sec:gaps} motivate a common reference structure for
information crossing a multi-robot execution boundary.  The framework does not
replace the contract theories surveyed in \cref{sec:related-work}.  An
implementation still needs an appropriate behavioral model when it claims
refinement, realizability, controller synthesis, or a global system property.

\subsection{Contracts on Producer--Consumer Edges}

Represent the deployed system as a directed graph.  A vertical edge
$e^{\mathrm v}_{r,i}:(r,i)\rightarrow(r,i+1)$ connects functions associated with
robot $r$, while a horizontal edge
$e^{\mathrm h}_{r,i\rightarrow s,j}:(r,i)\rightarrow(s,j)$ connects robots or
distributed components.  A horizontal edge is not merely another message path:
its consumer can act on behalf of a different robot and thereby change joint
feasibility, resource ownership, or a peer's safety margin.

For a producer--consumer edge $e=(p,c)$, define the reference contract
\begin{equation}
  C_e=(P_e,M_e;\,A_e,G_e,V_e,E_e,R_e),
  \label{eq:contract}
\end{equation}
where $P_e$ is a typed payload, $M_e$ is typed metadata, $A_e$ is the
consumer's assumption, and $G_e$ is the producer's guarantee.  $V_e$ contains
validators or monitor references, $E_e$ supplies supporting evidence such as a
solver result, proof, calibration report, monitor record, or test identifier,
and $R_e$ defines the response to rejection or unresolved validity.

Fields accompanying a message become contractual only when they participate in
explicit assumptions, guarantees, admission predicates, evidence requirements,
and recovery behavior.
\Cref{tab:coverage-obligations} maps each survey gap to the corresponding
metadata, validator, and recovery obligation.  This mapping is the bridge from
the literature synthesis to the proposed reference structure.

\begin{table}[!ht]
  \centering
  \caption{From survey gaps to edge-contract obligations}
  \label{tab:coverage-obligations}
  \scriptsize
  \setlength{\tabcolsep}{3pt}
  \begin{tabularx}{\textwidth}{@{}P{0.16\textwidth}P{0.32\textwidth}P{0.24\textwidth}Y@{}}
    \toprule
    Coverage gap & Required contract information & Required validator & Rejection or recovery \\
    \midrule
    Semantic and physical validity & Canonical task and resource identifiers; hard constraints; frame and units; model, operating region, and margins & Registry resolution and task, feasibility, or certificate check & Clarify, reject, repair the task, or replan \\
    Temporal validity & Event or forecast-target time; publication and availability observations; validity interval; clock domain and error bound & Freshness, target-time, deadline, and clock-comparability predicates & Request fresh state, shorten the horizon, degrade, or stop \\
    Uncertainty, provenance, and evidence & Uncertainty type, scope, correlation, calibration, propagation rule, model and solver identity, lineage, and evidence kind & Uncertainty-threshold, calibration, certificate, and provenance-policy checks & Re-estimate, re-solve, choose a robust alternative, or escalate \\
    Compatibility and lifecycle & Schema, protocol, implementation, dependencies, configuration, membership, resource ownership, and team epoch & Compatibility, supported-combination, epoch, ownership, and lifecycle-transition checks & Reconcile, downgrade, isolate, restart, or reallocate \\
    Intervention and authority & Producer and consumer authority; safety override; acknowledgement; completion evidence; recovery owner & Authority, acknowledgement, intervention-report, and completion predicates & Transfer authority, coordinate replanning, stop, or request human action \\
    \bottomrule
  \end{tabularx}
\end{table}


\Cref{fig:contract-lifecycle} previews how these obligations drive admission,
rejection, recovery, and the creation of new lineage.

\begin{figure}[!t]
  \centering
  \begin{tikzpicture}[
    x=1cm,y=1cm,
    box/.style={draw,rounded corners=2pt,minimum height=7mm,align=center,
      font=\footnotesize,inner sep=2.5pt},
    arrow/.style={-{Latex[length=2mm]},line width=0.7pt},
    tag/.style={font=\scriptsize,align=center,inner sep=1pt}
  ]
    \definecolor{ContractBlue}{HTML}{DCEBFA}
    \definecolor{ValidatorGreen}{HTML}{DDEFD8}
    \definecolor{RecoveryOrange}{HTML}{F9E4C8}
    \definecolor{RejectRed}{HTML}{F7D8D5}

    \node[box,fill=ContractBlue,minimum width=36mm] (producer) at (0,0)
      {Producer output};
    \node[box,fill=ContractBlue,minimum width=52mm] (contract) at (0,-1.05)
      {Interface contract $C_e$: payload,\\metadata, assumptions, evidence};
    \node[box,fill=ValidatorGreen,minimum width=64mm] (validators) at (0,-2.35)
      {Required validators $V_{e,q}$, $q\in\mathcal Q_e$: syntax, time,\\
       uncertainty, compatibility, $A_e$, $G_e$, provenance};

    \coordinate (branch) at (0,-3.25);
    \node[box,fill=ValidatorGreen,text width=27mm] (admit) at (-2.05,-4.05)
      {Admit to consumer};
    \node[box,fill=RejectRed,text width=30mm] (reject) at (1.95,-4.05)
      {Reject; record failed predicate and evidence};
    \node[box,fill=RecoveryOrange,text width=30mm] (recover) at (1.95,-5.25)
      {Invoke $R_e$: refresh, re-estimate, replan, degrade, or stop};
    \node[box,fill=ContractBlue,text width=30mm] (repair) at (1.95,-6.40)
      {New contract instance with new lineage};

    \draw[arrow] (producer) -- (contract);
    \draw[arrow] (contract) -- (validators);
    \draw[line width=0.7pt] (validators.south) -- (branch);
    \draw[arrow] (branch) -| (admit.north);
    \draw[arrow] (branch) -| (reject.north);
    \node[tag,anchor=south] at (-2.05,-3.22) {all required $=\true$};
    \node[tag,anchor=south] at (1.95,-3.22) {\false{} or required \unknown};
    \draw[arrow] (reject) -- (recover);
    \draw[arrow] (recover) -- (repair);
    \draw[arrow] (repair.east) -- ++(0.35,0) |- (contract.east);
  \end{tikzpicture}
  \caption{Admission and recovery on one interface.  A required check that
  returns \unknown{} blocks admission like \false.  A successful repair creates a
  new contract instance and provenance record while preserving the rejected
  claim.}
  \label{fig:lifecycle}
\end{figure}


Provenance can use the entity--activity--agent relations of W3C PROV
\cite{w3c2013prov}.  The framework keeps schema, protocol, and implementation
identity separate because syntactic compatibility does not establish behavioral
compatibility.

At an edge, local compatibility requires the producer guarantee to imply the
consumer assumption,
\begin{equation}
  G_e\Longrightarrow A_e.
  \label{eq:contract-composition}
\end{equation}
A consumer with several peer inputs must additionally check joint assumptions,
such as a common team epoch, consistent frames, unique resource ownership, and a
reconciliation rule after a partition.  Per-message checks cannot establish
these relations independently.  Equation~\eqref{eq:contract-composition} is
therefore a local condition, not proof that the composed team satisfies a global
mission property.  Such a proof requires defined input/output ownership,
feedback and temporal semantics, and an appropriate composition or refinement
calculus \cite{luckcuck2022compositional,mallozzi2022repair}.

\subsection{Admission and Recovery}

A syntactic validator checks decoding, required fields, types, dimensions,
units, frames, and schema compatibility.  A semantic validator may call a
solver, runtime monitor, perception check, calibrated statistical test, or
certificate verifier.  Its result lies in
$\{\textsc{true},\textsc{false},\textsc{unknown}\}$.  Checks designated as
admission requirements accept only \textsc{true}; an advisory check may retain
\textsc{unknown} only when it is explicitly excluded from admission.

Let $\mathcal Q_e$ contain the required syntactic, temporal, uncertainty,
compatibility, assumption, guarantee, and provenance checks for edge $e$.  The
admission decision is
\begin{equation}
  \operatorname{Admit}_e(t)=\textsc{true}
  \quad\Longleftrightarrow\quad
  \bigwedge_{q\in\mathcal Q_e}
  \left[V_{e,q}(P_e,M_e,E_e,t)=\textsc{true}\right].
  \label{eq:contract-admission}
\end{equation}
Equation~\eqref{eq:contract-admission} is what turns descriptive fields into an
enforced boundary: every required predicate participates in the execution
decision.

At minimum, admission checks the relevant identity and frame, time and validity,
uncertainty threshold, supported implementation and configuration, consumer
assumptions, producer evidence, and provenance policy.  A \textsc{false} or
required \textsc{unknown} invokes $R_e$: request clarification or fresh data,
re-estimate, replan, degrade, isolate, stop, or escalate.  The rejection record
names the failed predicate, time, evidence, and responsible boundary.  A repaired
result receives new lineage and does not silently replace the rejected claim.

\subsection{Failure Attribution}

The same physical outcome can have several origins.  A collision may begin with
an incorrect semantic relation, a biased prediction, a missing physical
constraint, stale peer state, or a hardware fault.  Conversely, downstream
checks can turn the same originating fault into a rejection, degraded mission,
or safe stop.  The taxonomy in \cref{tab:failure-propagation} therefore records
origin, contract concern, first observation, effect, and recovery authority
separately.

\FloatBarrier
\begin{table}[!t]
  \centering
  \caption{Proposed cross-layer failure taxonomy.  First observation denotes the earliest boundary where fault evidence can become available}
  \label{tab:failure-propagation}
  \scriptsize
  \begin{tabularx}{\textwidth}{@{}P{0.10\textwidth}P{0.14\textwidth}P{0.14\textwidth}Y P{0.15\textwidth}P{0.14\textwidth}@{}}
    \toprule
    Origin & Local fault & Contract concern & First observation & Possible effect & Recovery authority \\
    \midrule
    Perception / estimation & Misassociation, bias, or clock error & Payload, frame, time, uncertainty & Estimator innovation, semantic grounding, or controller residual & Wrong referent, localization drift, or unsafe state & Re-estimation, sensor fallback, or stop \\
    Semantic planning & Invented relation, capability, or priority & Intent, assumptions, guarantees & Task-model validator or downstream feasibility check & Wrong or impossible assignment; repeated replanning & Clarify, repair, or replan \\
    World model & Biased or overconfident rollout & Horizon, uncertainty type, validity evidence & Online correction, twin residual, or physics validator & Coordination failure or poor recovery & Re-estimate, shorten horizon, or choose robust plan \\
    Physical model & Missing constraint or inaccurate contact & Feasibility claim and supporting solver evidence & Optimizer, controller residual, or contact event & Slip, saturation, collision, or damage & Re-optimize, fallback controller, or stop \\
    Runtime / network & Delay, loss, partition, version, or lifecycle fault & Age, clock, protocol, implementation, assumptions & Consumer monitor, missed update, exception, or teardown & Inconsistent team state, abort, deadlock, or collision & Degrade, reconnect, reallocate, or restart \\
    Safety function & Local or unreported intervention & Authority, intervention evidence, recovery & Safety monitor or upstream acknowledgement & Safe local motion but failed mission or new conflict & Coordinated replan or human escalation \\
    Hardware & Sensor, actuator, energy, or compute fault & Capability and health guarantee & Health monitor, tracking residual, or peer timeout & Task abandonment or unsafe peer compensation & Isolate robot and dynamically reallocate \\
    \bottomrule
  \end{tabularx}
\end{table}


\subsection{Worked Command Trace}

Consider the instruction: \emph{``Use two robots to inspect sector B while
keeping the corridor clear.''}  \Cref{tab:command-trace} follows the command
from resolved roles and constraints to trajectories, runtime state, safety
intervention, and completion evidence.  Each row includes the information
required by the next validator as well as the nominal output.  It is a reference
representation, not a universal ROS message definition; thresholds and recovery
policies remain scenario-specific.

\FloatBarrier
\begin{table}[t]
  \centering
  \caption{Running command trace; all thresholds are scenario-specific contract parameters}
  \label{tab:command-trace}
  \scriptsize
  \begin{tabularx}{\textwidth}{@{}P{0.13\textwidth}P{0.27\textwidth}Y Y@{}}
    \toprule
    Function & Payload $P_e$ & Assumption, guarantee, and metadata & Validator and rejection $V_e,R_e$ \\
    \midrule
    Semantic task formation & Roles, sector identifier, inspection goal, and corridor-clear constraint & $A_e$: referents resolve and two capable robots exist; $G_e$: confirmed hard constraints remain explicit & Resolve against map and capability registries; clarify or reject an unresolved referent \\
    Predictive state & Occupancy and joint trajectories over horizon $H$ & $M_e$: event and target times, horizon, uncertainty semantics, and other-agent policy identity & Reject or shorten $H$ if age exceeds $\bar\Delta$ or uncertainty is not below its declared limit $\bar U$ \\
    Physical grounding & Reachable viewpoints and time-parameterized paths & $G_e$: separation is at least $d_{\min}$ and actuator limits hold; $E_e$: model, solver status, and minimum margins & Re-solve or return infeasible when solver status is unacceptable or a required margin is negative \\
    Distributed runtime & Assignments, paths, peer state, and acknowledgements & $M_e$: frame, units, schema/protocol/implementation identities, team epoch, ownership, and clock-error bound & Drop expired or wrong-epoch data; reconcile, degrade, or replan when age, clock, ownership, or compatibility checks fail \\
    Runtime safety & Admissible command or intervention record & $G_e$: conditional local-safety claim; $E_e$: active constraint, model/state/update assumptions, and intervention cause & Apply a feasible override and publish its cause; stop or escalate if the safety program is infeasible \\
    Physical robots & Execution state, health, resource release, and completion evidence & Capability and health guarantee; ownership and completion status & Isolate an unhealthy robot, reallocate its task, or request human action \\
    \bottomrule
  \end{tabularx}
\end{table}

